\chapter{Research Design and Methodology}
\label{ch:methodology}

\section{Research Design and Mode}
\label{sec:research-mode}

This dissertation adopts an \textbf{empirical systems engineering} paradigm~\cite{Wohlin2012,Runeson2009} to evaluate the architectural, performance, and runtime consequences of migrating a Spring Boot Model-View-Controller (MVC) system to Java~21 Virtual Threads (Project Loom). Rather than treating runtime migration as a trivial configuration toggle, this research models the transition as an architectural intervention within a multi-tier concurrency and queueing network.

The empirical investigation is structured around a controlled, instrumented Spring Boot MVC benchmark application (\texttt{virtual-threads-demo}). In empirical software engineering, research designs must balance internal validity (the degree to which observed outcomes are directly attributable to experimental treatments) against external validity (the generalizability of conclusions to industrial systems):
\begin{itemize}
  \item \textbf{Limits of Synthetic Micro-benchmarks:} Instruction-level benchmarks implemented via the Java Microbenchmark Harness (JMH) excel at isolating CPU cache and branch behavior, but completely omit framework-level interactions: HTTP socket multiplexing, servlet container request dispatching via \texttt{newVirtualThreadPerTaskExecutor()}, connection pool checkout arbitration, and request-scoped \texttt{ThreadLocal} lifecycles.
  \item \textbf{Limits of Production Deployment Testing:} Full-scale production systems introduce severe confounding environmental variables---wide-area network jitter, database optimizer variance, multi-tenant CPU throttling, and background operating system daemons---that obscure the low-level scheduling and synchronization mechanics of the JVM.
  \item \textbf{The Representative Testbed Strategy:} A dedicated, representative framework artifact provides an optimal balance: it preserves the authentic request lifecycle of Spring MVC, Apache Tomcat, and HikariCP while establishing strict experimental controls over carrier concurrency, workload generation, lock boundaries, and flight recorder telemetry.
\end{itemize}

\section{Question-to-Method Operationalisation}
\label{sec:rq-operationalisation}

To transition the two research questions (RQ1 and RQ2) into empirically verifiable hypotheses, each question is operationalized around explicit treatments, controlled baselines, dependent metrics, and decision rules. Table~\ref{tab:rq-method} details this operational mapping.

\begin{table}[H]
  \centering
  \small
  \caption{Research-question-to-method operationalisation matrix.}
  \label{tab:rq-method}
  \begin{tabularx}{\textwidth}{@{}lp{0.22\textwidth}p{0.18\textwidth}p{0.18\textwidth}X@{}}
    \toprule
    \textbf{RQ} & \textbf{Experimental Treatment} & \textbf{Controlled Baseline} & \textbf{Dependent Metrics} & \textbf{Operational Decision Rule} \\
    \midrule
    RQ1 & Virtual Threads + \texttt{synchronized} & Platform Threads + \texttt{synchronized} & Throughput ($X$), $p50$, $p95$, $p99$ latency, JFR pins & Carrier pinning confirmed if JFR records events $\ge 20\text{ ms}$; carrier starvation confirmed if throughput collapses under \texttt{parallelism=1}. \\
    \addlinespace
    RQ1 & Virtual Threads + \texttt{ReentrantLock} & Virtual Threads + \texttt{synchronized} & Throughput ($X$), latency distributions, $\sigma$ dispersion & Unpinning confirmed if JFR events $= 0$; Amdahl serialization confirmed if throughput remains within $\pm 5\%$ of service capacity $\mu = 1/D = 2.0\text{ req/s}$. \\
    \addlinespace
    RQ1 & AQS non-fair lock dispersion & Virtual Threads + \texttt{synchronized} & Latency standard deviation ($\sigma$) and range ($\Delta$) & AQS non-fair lock trial-level dispersion confirmed if \texttt{ReentrantLock} standard deviation exceeds \texttt{synchronized} by $> 5\times$ across process repetitions. \\
    \addlinespace
    RQ1 & Java 24 JEP 491 runtime & Java 21 LTS Virtual Threads + \texttt{synchronized} & JFR pinning count, throughput ($X$) & JEP 491 monitor unpinning confirmed if JFR pinning events drop to strictly 0 while throughput remains bounded by Amdahl ceiling $1/D$. \\
    \addlinespace
    RQ2 & Native JNI POSIX \texttt{nanosleep()} & Pure Java managed \texttt{Thread.sleep()} & Throughput ($X$), JFR pinning count, process \%CPU & Carrier capture confirmed if throughput drops to serialized Amdahl rate $1/D$ while JFR pinning remains 0 and CPU remains negligible (non-spinning kernel wait). \\
    \addlinespace
    RQ2 & Persistence holding (Stratum~8) & Pre-acquisition delay (Treatment A) \& Pool 10 & Throughput ($X$), HikariCP active/pending queues & Pool saturation confirmed if single-connection pool caps throughput to $X \le 1/E[T_{\text{hold}}]$ with queuing, while expanded pool scales linearly up to concurrency bounds with 0 JFR pins. \\
    \bottomrule
  \end{tabularx}
\end{table}

\section{Materials, Data, Tools, and Provenance}
\label{sec:materials-provenance}

The experimental testbed comprises open-source software libraries, standard compilation toolchains, and dedicated bare-metal hardware:
\begin{itemize}
  \item \textbf{Hardware Platform:} Dedicated bare-metal Apple M1 processor (8 physical execution cores: 4 Firestorm performance cores, 4 Icestorm efficiency cores), 8~GB unified LPDDR4X memory, running Darwin 25.2.0 (macOS 15.2, arm64) for Strata 1--7 and Darwin 27 for the November Stratum~8 runs.
  \item \textbf{Software Stack:} Java 21 LTS (Azul Zulu 21.0.8 for Strata 1--7 and Zulu 21.0.12.1 for Stratum~8, OpenJDK HotSpot 64-bit Server VM), OpenJDK 24.0.2 (build 24.0.2+12), Spring Boot 3.5.12, embedded Apache Tomcat 10.1.52, HikariCP 6.3.3, and PostgreSQL 42.7.10 JDBC driver.
  \item \textbf{Workload Generator:} Grafana k6~\cite{GrafanaK6} compiled for Darwin arm64 (v1.7.1 in Strata 1--7; v2.3.0 in Stratum~8).
  \item \textbf{Telemetry Daemons:} JDK Flight Recorder (JFR) loaded via HotSpot \path{profile.jfc}, a BSD \texttt{ps} polling daemon sampling process CPU utilization at 1.0-second intervals, and Spring Boot Actuator HikariCP metric polling (\path{/actuator/metrics/hikaricp.connections.*}) at 1.0-second intervals in Stratum~8.
\end{itemize}

\subsection{Repository Evidence Strata Census}
To preserve provenance traceability and prevent conflating exploratory runs with primary benchmarks, the corpus of 125 physical trials (+1 schema template) is partitioned into eight non-overlapping strata, as detailed in Table~\ref{tab:strata-census}. Throughout this chapter and Chapters~\ref{ch:results}--\ref{ch:discussion}, parenthetical identifiers such as \texttt{CLM-RQ1-01} or \texttt{CLM-METH-02} trace directly to the repository's frozen claim-to-evidence ledger (\path{docs/evidence-freeze/claim-evidence-ledger.md}).

\begin{table}[H]
  \centering
  \caption{Master repository execution census across eight evidence strata.}
  \label{tab:strata-census}
  \begin{tabularx}{\textwidth}{@{}p{0.14\textwidth}p{0.26\textwidth}cp{0.16\textwidth}X@{}}
    \toprule
    \textbf{Stratum} & \textbf{Focus Description} & \textbf{Trials} & \textbf{Admitted Status} & \textbf{Epistemic Role} \\
    \midrule
    Stratum 1 & Java 21 Primary Synchronization Matrix & 12 & Admitted (Primary) & Mutex treatments (\texttt{synchronized}, \texttt{ReentrantLock}, no-lock) \\
    Stratum 2 & Java 24/25 Runtime Evolution Matrix & 10 & Admitted (Supporting) & JEP 491 object monitor unpinning verification \\
    Stratum 3 & Multi-JDK Native JNI Boundary Matrix & 27 & Admitted (Primary) & $3 \times 3 \times 3$ factorial evaluation of JNI POSIX sleep \\
    Stratum 4 & Supporting Architectural Controls & 18 & Admitted (Supporting) & Endpoint controls: \texttt{/orders}, \texttt{/products}, \texttt{/cpu} \\
    Stratum 5 & 8-Carrier Historical Baseline Matrix & 12 & Admitted (Baseline) & Unconstrained 8-carrier Apple M1 baseline \\
    Stratum 6 & Exploratory Pilots \& Diagnostic Runs & 27 & Quarantined & Preliminary lock/semaphore exploratory runs \\
    Stratum 7 & Construct-Flawed / Defective Runs & 13 & Quarantined & Pre-acquisition sleep flaws, unmanifested runs \\
    Stratum 8 & Two-Treatment Persistence Ablation & 6 & Admitted (Primary) & HikariCP pool size 1 vs.\ 10 connection-holding ablation \\
    Template & Run Manifest Reference Schema & 1 & Metadata Only & Reference JSON schema specification \\
    \midrule
    \textbf{Total} & \textbf{Complete Repository Inventory} & \textbf{126} & \textbf{543 Artifacts} & \textbf{Cryptographically SHA-256 verified} \\
    \bottomrule
  \end{tabularx}
\end{table}

\section{Workload Protocol and Closed-Loop Queueing Mechanics}
\label{sec:workload-protocol}

Benchmark load is generated using a staged closed-loop concurrency profile executed over 55.0 seconds per trial:
\begin{enumerate}
  \item \textbf{Stage 1 (Ramp-Up, $t \in [0\text{ s}, 15\text{ s}]$):} Concurrency scales linearly from 0 to $N$ Virtual Users (VUs), warming JIT compilation and populating thread pools.
  \item \textbf{Stage 2 (Steady-State Hold, $t \in [15\text{ s}, 45\text{ s}]$):} Load is maintained at constant concurrency $N = 10\text{ VUs}$ (Strata 1--5) or $N = 50\text{ VUs}$ (Stratum 8) for exactly 30.0 seconds.
  \item \textbf{Stage 3 (Ramp-Down, $t \in [45\text{ s}, 55\text{ s}]$):} Concurrency ramps down linearly from $N$ to 0.
\end{enumerate}

\subsection{Closed-Loop vs. Open-Loop Mechanics}
In accordance with Schroeder et al.~\cite{Schroeder2006}, selecting a closed-loop workload model enforces client-side self-throttling. Each virtual user issues an HTTP request, awaits the response, executes a constant think-time delay ($Z = 100\text{ ms} = 0.100\text{ s}$), and then submits the subsequent request.

By Little's Law for interactive closed-loop systems~\cite{Little1961}:
\begin{equation}
  \label{eq:interactive-law}
  N = X \times (R + Z) \iff R = \frac{N}{X} - Z
\end{equation}
where $N$ is the active virtual user population, $X$ is throughput ($\text{req/s}$), $R$ is residence/response time ($\text{s}$), and $Z$ is think time ($\text{s}$).

For a serialized critical section with duration $D = 0.500\text{ s}$, the theoretical service capacity ceiling is:
\begin{equation}
  \mu = \frac{1}{D} = \frac{1}{0.500\text{ s}} = 2.000\text{ req/s}
\end{equation}
Under $N = 10\text{ VUs}$ and $Z = 0.100\text{ s}$, the theoretical saturated cycle time is:
\begin{equation}
  R + Z = \frac{N}{\mu} = \frac{10}{2.000\text{ req/s}} = 5.000\text{ s} \implies R = 4.900\text{ s}
\end{equation}
Any experimental observation deviating from this theoretical expectation indicates either additional service overhead or queueing departure from stationary equilibrium.

\section{Dual-Window Steady-State Analysis Standard}
\label{sec:dual-window-standard}

A central methodological contribution of this dissertation is the formalization of the \textbf{Dual-Window Reporting Standard} (addressing \texttt{CLM-METH-02}):
\begin{itemize}
  \item \textbf{Scenario-Start Clock Anchoring ($t_0$):} To prevent initiation delays from misaligning metric time windows, the temporal origin $t_0$ is anchored at the earliest request initiation timestamp across all completed client requests ($t_{\text{init}} = t_{\text{completion}} - \text{duration}$).
  \item \textbf{Whole-Run Window:} Computes metrics across all completed requests over the observed test duration ($t_{\text{max}} - t_0$).
  \item \textbf{Steady-State Hold Window ($t \in [15\text{ s}, 45\text{ s}]$):} Isolates requests whose completion timestamps fall strictly within $t_{\text{completion}} - t_0 \in [15.0\text{ s}, 45.0\text{ s}]$, evaluated against the fixed 30.0-second hold denominator.
\end{itemize}

During the 15-second ramp-up stage, the server queue is sparsely populated, yielding transient response times between $500\text{ ms}$ and $3\,900\text{ ms}$. In ramp-sensitive workloads (such as \texttt{platform-synchronized}), whole-run median latency was $3\,917.56\text{ ms}$ compared to steady-state hold median latency of $4\,931.24\text{ ms}$. Aggregating across the whole run thus introduces a \textbf{20.56\% shortfall} relative to the true steady-state queueing baseline (or conversely, steady-state hold reflects a \textbf{25.87\% increase} over whole-run aggregation). Reporting both windows side-by-side guarantees complete transparency and eliminates warm-up bias~\cite{Traini2022}.

\section{Telemetry Pipeline and Correctness Invariants}
\label{sec:telemetry-correctness}

The experimental infrastructure coordinates three independent telemetry layers:
\begin{enumerate}
  \item \textbf{Client Stream (k6):} Raw JSON lines streams recording microsecond-accurate request start timestamps, durations, and HTTP response codes.
  \item \textbf{JVM Runtime Stream (JFR):} Continuous flight recordings capturing \texttt{jdk.VirtualThreadPinned} events. Crucially, HotSpot filters pinning events shorter than $20\text{ ms}$ (\texttt{threshold=20ms}). Thus, zero recorded events confirms the absence of pinning events $\ge 20\text{ ms}$, but does not preclude sub-20ms micro-pins (\texttt{CLM-METH-04}).
  \item \textbf{Host Telemetry Stream (OS):} Periodic process sampling recording overall JVM CPU utilization ($\%cpu$). During JNI POSIX sleep, process CPU averages approximately $1\%$, proving that the carrier thread is suspended in a kernel wait state rather than busy-spinning (\texttt{CLM-RQ2-02}).
\end{enumerate}

Every trial admitted into the primary corpus must satisfy six correctness gates: HTTP 200 status on 100\% of requests, sequential consistency in transactional endpoints, clean zero-exit process termination, driver carrier-safety verification (\texttt{pgjdbc >= 42.6.0}), cryptographic SHA-256 manifest matching, and membership in pre-registered evidence strata.

\section{Pre-Execution Validity Controls}
\label{sec:validity-controls}

To protect internal, measurement, and statistical validity during data collection, five experimental controls were enforced prior to analysis (evaluated alongside external and construct threats in the master threats-to-validity matrix, Table~\ref{tab:master-validity} in Section~\ref{sec:master-validity}):
\begin{itemize}
  \item \textbf{Single-Carrier Isolation (\texttt{parallelism=1}):} Mechanism-isolation trials constrain the JVM \texttt{ForkJoinPool} to one carrier worker so monitor pinning cannot hide behind spare carriers, while Stratum~5 captures unconstrained 8-carrier behavior.
  \item \textbf{Closed-Loop Stationary Equilibrium:} Fixed VU counts ($N=10$ or $50$) with $Z=100\text{ ms}$ think time prevent client packet drops and permit exact Little's Law verification ($R = N/X - Z$).
  \item \textbf{Dual-Window Reporting \& Descriptive $N=3$ Bounds:} Reporting whole-run alongside hold-window ($15\text{--}45\text{ s}$) metrics eliminates warm-up ramp distortion, while $N=3$ independent process runs are reported strictly as descriptive bounds ($\sigma, \Delta$) without inferential $p$-values or ANOVA.
  \item \textbf{Telemetry Threshold Qualification:} Because HotSpot \path{profile.jfc} filters \texttt{jdk.VirtualThreadPinned} events below $20\text{ ms}$, zero recorded events indicate absence of pins $\ge 20\text{ ms}$ rather than absence of sub-20\,ms micro-pins.
  \item \textbf{Quarantine \& Inventory Integrity:} Construct-flawed or uninstrumented runs (pre-acquisition sleep in Stratum~7; Docker PostgreSQL runs lacking JFR) are quarantined in Strata~6--7, and all 543 physical artifacts are tracked against SHA-256 checksums.
\end{itemize}

\section{Academic Ethics and AI Disclosure}
\label{sec:ethics-ai-methodology}

AI assistance is documented under Dean's Directive~1/2026~\cite{elteikAiPolicy} in Appendix~\ref{app:ethics-ai} (Table~\ref{tab:ai-usage-breakdown}). Zero empirical benchmark measurements, latency figures, or flight recordings were generated by artificial intelligence; all empirical findings originate from physical hardware executions verified via cryptographic hashing and independent recalculation.
